\section*{PROCEDIMIENTOS PARA LA RECUPERACIÓN DE FUENTES}

Una vez que el problema se presenta y la fuente está fuera de control, las acciones que se sigan deberán programarse de tal manera que el personal involucrado reciba la mínima dosis posible. Debemos recordar que en la mayoría de los accidentes de radiografía industrial (no incluyendo la pérdida del contenedor y de la fuente) no hay de por medio vidas que pudieran estar en peligro, por lo que hay tiempo suficiente para programar las actividades que deban realizarse. Actuar precipitadamente por lo general da malos resultados.

A continuación describiremos algunos procedimientos para evitar que la pérdida de control de la fuente sea de lamentables consecuencias y de cómo recuperar el control, colocando y asegurando la fuente en el contenedor.
